\documentclass{article}

\usepackage{changepage}
\usepackage{xcolor}
\usepackage{amsmath}
\usepackage{amssymb}
\usepackage{tikz}
\usepackage{tcolorbox}

\definecolor{blue}{RGB}{0, 166, 223}
\definecolor{lightblue}{RGB}{220, 240, 247}

\newcommand{\header}[1]{
	\vspace{0.25cm}
		\begin{adjustwidth}{0cm}{0cm}
			\textcolor{blue}{\textbf{#1}}
		\end{adjustwidth}
	\vspace{0.25cm}
}

\newcommand{\question}[2]{
	\vspace{0.125cm}
		\begin{adjustwidth}{0.25cm}{0cm}
			\textbf{#1} #2
		\end{adjustwidth}
	\vspace{0.125cm}
}

\newcommand{\container}[1]{
	\vspace{0.25cm}
		\begin{adjustwidth}{0cm}{0cm}
			#1
		\end{adjustwidth}
	\vspace{0.25cm}
}

\newcommand{\solution}[1]{
	\vspace{0.25cm}
		\begin{adjustwidth}{0cm}{0cm}
			\textbf{Solution:} #1
		\end{adjustwidth}
	\vspace{0.25cm}
}

\newcommand{\wrapper}[1]{
	\vspace{0.25cm}
		\begin{adjustwidth}{0cm}{0cm}
			#1
		\end{adjustwidth}
	\vspace{0cm}
}

\newcommand{\calc}[1]{
	\[
		#1
	\]
}

\begin{document}
	
	\header{5.2.2}

	\container{
		For each positive integer \emph{n}, let \emph{P(n)} be
		the formula

		\calc{1^{2} + 2^{2} + \text{...} + n^{2} = \frac{n(n + 1)(2n + 1)}{6}}

		\question{a.}{Write \emph{P(1)}. Is \emph{P(1)} true?}
		\question{b.}{Write \emph{P(k)}.}
		\question{c.}{Write \emph{P(k + 1)}.}
		\question{d.}{
			In a proof by mathematical induction that the
			formula holds for every integer \(n \geq 1\),
			what must be shown in the inductive step?
		}
	}

	\container {
		Let \emph{P(n)} be

		\calc{
			\sum_{i=1}^{n} i^{2} = \frac{n(n + 1)(2n + 1)}{6}
		}
	}

	\container {
		\textbf{a.}

		\wrapper{For the left hand side:}

		\calc{
			P(1) = \sum_{i=1}^{1} i^{2} = (1)^{2} = 1^{2} = \boxed{1}
		}

		\wrapper{For the right hand side:}

		\calc{
			P(1) = \frac{1(1 + 1)(2(1) + 1)}{6} = \frac{2 \times 3}{6} = \frac{6}{6} = \boxed{1}
		}
	}

	\solution{\emph{P(1)} is true.}

	\container {
		\textbf{b.}

		\wrapper{For the left hand side:}

		\calc{
			P(k) = \sum_{i=1}^{k} i^{2} = \boxed{1^{2} + 2^{2} + \text{...} + k^{2}}
		}

		\wrapper{For the right hand side:}

		\calc{
			P(k) = \boxed{\frac{k(k + 1)(2k + 1)}{6}}
		}
	}

	\solution{\(P(k) = 1^{2} + 2^{2} + ... + k^{2} = \frac{k(k + 1)(2k + 1)}{6}\)}

	\container {
		\textbf{c.}
		
		\wrapper{For the left hand side:}

		\calc{
			\sum_{i=1}^{k + 1} i^{2} = 1^{2} + 2^{2} + \text{...} + k^{2} + (k + 1)^{2}
		}

		\wrapper{Let \(\frac{k(k + 1)(2k + 1)}{6} = 1^{2} + 2^{2} + \text{...} + k^{2}\).}

		\calc{
			\sum_{i=1}^{k + 1} i^{2} = \frac{k(k + 1)(2k + 1)}{6} + (k + 1)^{2}
		}

		\wrapper{For the right hand side:}

		\calc{
			\frac{(k + 1)((k + 1) + 1)(2(k + 1) + 1)}{6} = \boxed{\frac{(k + 1)(k + 2)(2k + 3)}{6}}
		}
	}

	\solution{}

	\container {
		\textbf{d.}
		
		\wrapper{
			In a proof by mathematical induction that the
			formula holds for every integer \(n \geq 1\),
			we suppose that the formula is true for a particular
			but arbitrarily chosen integer \emph{k}, and if true,
			it must also be true for \emph{k + 1}, for both
			the left-hand side and the right-hand side.
		}
	}

	\solution{See above.}
\end{document}
